Repository-level repair systems emit many logged events before a verifier pass. A rate over one event is easy to misread when the report omits the event's operational meaning, opportunity budget, support unit, or verifier boundary. We introduce a claim-indexed measurement model that binds each repair statistic to five coordinates: construct, opportunity design, unit of support, endpoint interpretation, and instrumentation provenance. We instantiate the model in a retrospective audit of \TotalRows{} executions: two 30-task DeepSeek strata with disjoint task IDs and a 30-task GLM boundary stratum that reuses the clean task surface. A field-integrity audit identifies \ContradictionRows{} raw GLM action positives that conflict with the patch-application evidence required by the governed-action construct. The harmonized table contains no verifier pass outside the governed-action state, while conditional endpoint yield is \CleanPasses{}/\CleanActions{} (\CleanYield\%) and \ExtensionPasses{}/\ExtensionActions{} (\ExtensionYield\%) in the two DeepSeek strata. Under eight retained opportunities per fixed cell, governed-action discovery reaches \CleanActionDiscoveryKEight\% and \ExtensionActionDiscoveryKEight\%; verifier-pass discovery reaches \CleanPassDiscoveryKEight\% and \ExtensionPassDiscoveryKEight\%. The \DeepSeekPasses{} passing rows reduce to seven positive cells, three tasks, and two task families. Endpoint controls produce the expected result for all six gold runs, six no-op runs, and nine reference-derived mutants. These findings show that repair evidence cannot be reduced to one success scalar. Field semantics, opportunity, support, verifier scope, and instrumentation determine what a reported number measures and which comparisons it can sustain.
